% Første linje er dokumentklassen. Brug *altid* memoir!
\documentclass[a4paper,oneside,article]{memoir}

% Herunder indlæses forskellige pakker
\usepackage[utf8]{inputenc}  % Korrekt håndtering af æ, ø og å
\usepackage[T1]{fontenc}  % Korrekt håndtering af æ, ø og å
\usepackage{microtype}  % Typografisk magi! Giver bl.a. pænere orddeling
\usepackage{graphicx}  % Gør det muligt at indsætte billeder
\usepackage{amsmath}  % Giver adgang til uundværlige matematikting
\usepackage{mathtools}  % Retter ting i ams og giver ekstra muligheder
\usepackage[danish]{babel}  % Danske betegnelser og orddeling
\usepackage{lmodern}
\usepackage{alphabeta}
\usepackage{float} %Mere kontrol over billede-placering
\renewcommand{\danishhyphenmins}{22}  % Bedre dansk orddeling
\usepackage[exponent-product=\ensuremath{\cdot}]{siunitx} % Sientific notation for forskellige ting, med en regel om at den skal bruge cdot istedet for times. 
\usepackage{derivative} % Til at skrive differentialligninger 
\usepackage{xcolor} % For at få flere tekst farver 
\usepackage[colorlinks=true]{hyperref} % For at få refrences man kan trykke på, slet [] hvis man vil have kasser i stedet for farvet tekst.
\usepackage{pdfpages} % Til at indsætte andre pdf'er
\usepackage{subfig} % For at have to billeder ved siden af hinanden
\usepackage{unicode-math}

% Pænere toc indents 
\usepackage{tocloft}
\cftsetindents{section}{1em}{2em}
\cftsetindents{subsection}{4em}{0em}


% Indholdet i dit dokument starter her!
\begin{document}
\author{Mikkel Moth Billing \\ 202307989}
\title{Videregående Mekanik Formelsamling}
\date{\today}
\maketitle   % Indsætter overskriften
% Table of Contents indstillinger
\setcounter{tocdepth}{2}
\setcounter{secnumdepth}{1}
\hypersetup{linkcolor=black} % Laver tableofcontents sort.  
\tableofcontents %Indsætter ToC
\hypersetup{linkcolor=red} % og så alle andre hyperrefs røde. 

%Custom Commands
%Til opgavebeskrivelser
\newcommand{\opg}[1]{\textcolor{gray}{\emph{#1}}\plainbreak{1}} 

\newcommand{\ti}[1]{\cdot 10^{#1}} % Let: gange 10^x
\newcommand{\e}[1]{\emph{e}^{#1}} % Pænere e^x
\newcommand{\vm}[2]{\begin{pmatrix} #1 \\ #2 \end{pmatrix}}
\newcommand{\mat}[1]{\begin{pmatrix} #1 \end{pmatrix}}

% Lette paranteser 
\newcommand{\br}[2][1]{%
    \ifcase#1
    \or \left( #2 \right) %1 eller ingenting
    \or \left[ #2 \right] %2
    \or \left\langle #2 \right\rangle %3
    \or \left\{ #2 \right\} %4
    \or \left| #2 \right| %5
    \fi
}
%Til at sætter bileder ind i store dokumenter
\newcommand{\bil}[1]{
\begin{figure}[H]
    \centering
    \includegraphics[width=0.75\linewidth]{Billeder/#1.png}
    \label{fig:#1}
\end{figure}
}

%Til at sætte pdf'er ind i store dokumenter (kun pdf'er på mere end 1 sider)
\newcommand{\pdf}[4]{
\includepdf[pages=#3,scale=.85,pagecommand={\subsection{#2 - \ref{ch:opg}}  \label{asc:#2}}]{pdf/#1.pdf}
\includepdf[pages=\the\numexpr 1 + #3\relax-#4, scale=.85]{pdf/#1.pdf}
}

\renewcommand{\vec}[1]{\textbf{#1}}

% Lav ny pagestyle (kun odd fordi enkeltsidet)
\makepagestyle{Title}
\pagestyle{Title}
%\makeoddhead{Title}{\footnotesize\color{gray}{\theauthor}}{}{\textcolor{gray}{\thedate}}
\makeoddfoot{Title}{}{\footnotesize\color{gray}{\thepage{} / \thelastpage}}{}
% Også fod på siden med \maketitle
\makeoddfoot{plain}{}{\footnotesize\color{gray}{\thepage{} / \thelastpage}}{}
\newpage

%Skriv
\chapter{Basics}
\section{Newton}
\begin{align*}
    \vec{F} = m\vec{a} = m\ddot{\vec{r}}=\dot{\vec{p}}
\end{align*}
Hvis der ikke er en udefrapåvirkende kraft er impulsen bevaret. 

\section{Trigonometri}
\begin{align*}
    &\sin(\theta \pm \phi) = \sin \theta \cos \phi \pm \cos \theta \sin \phi \\
    &\cos(\theta \pm \phi) = \cos \theta \cos \phi \mp \sin \theta \sin \phi \\
    &\cos \theta \cos \phi = \frac{1}{2} \left[\cos(\theta + \phi) + \cos(\theta - \phi)\right] \\
    &\sin \theta \sin \phi = \frac{1}{2} \left[\cos(\theta - \phi) - \cos(\theta + \phi)\right] \\
    &\sin \theta \cos \phi = \frac{1}{2} [\sin(\theta + \phi) + \sin(\theta - \phi)] \\
    &\cos^2 \theta = \frac{1}{2} [1 + \cos 2\theta] \\
    &\sin^2 \theta = \frac{1}{2} [1 - \cos 2\theta] \\
    &\cos \theta + \cos \phi = 2 \cos \left( \frac{\theta + \phi}{2} \right) \cos \left( \frac{\theta - \phi}{2} \right) \\
    &\cos \theta - \cos \phi = -2 \sin \left( \frac{\theta + \phi}{2} \right) \sin \left( \frac{\theta - \phi}{2} \right) \\
    &\sin \theta + \sin \phi = 2 \sin \left( \frac{\theta + \phi}{2} \right) \cos \left( \frac{\theta - \phi}{2} \right) \\
    &\sin \theta - \sin \phi = 2 \cos \left( \frac{\theta + \phi}{2} \right) \sin \left( \frac{\theta - \phi}{2} \right) \\
    &\cos^2 \theta + \sin^2 \theta = 1\\
    &\sec^2 \theta - \tan^2 \theta = 1\\
    &e^{i \theta} = \cos \theta + i \sin \theta \quad \text{[Euler's relation]} \\
    &\cos \theta = \frac{1}{2} \left( e^{i \theta} + e^{-i \theta} \right) \\
    &\sin \theta = \frac{1}{2i} \left( e^{i \theta} - e^{-i \theta} \right)
\end{align*}

\section{Series Expansions}
\begin{align*}
    f(z) = f(a) + & f'(a)(z - a) + \frac{1}{2!}f''(a)(z - a)^2 + \frac{1}{3!}f'''(a)(z - a)^3 + \dots \quad \text{[Taylor's series]} \\
    & e^z = 1 + z + \frac{1}{2!}z^2 + \frac{1}{3!}z^3 + \dots \\
    & \cos z = 1 - \frac{1}{2!}z^2 + \frac{1}{4!}z^4 - \dots \\
    & \cosh z = 1 + \frac{1}{2!}z^2 + \frac{1}{4!}z^4 + \dots \\
    & \tan z = z + \frac{1}{3}z^3 + \frac{2}{15}z^5 + \dots \quad [|z| < \pi/2] \\
    & (1 + z)^n = 1 + nz + \frac{n(n - 1)}{2!}z^2 + \dots \quad [|z| < 1] \quad \text{[binomial series]} \\
    & \ln(1 + z) = z - \frac{1}{2}z^2 + \frac{1}{3}z^3 - \dots \quad [|z| < 1] \\
    & \sin z = z - \frac{1}{3!}z^3 + \frac{1}{5!}z^5 - \dots \\
    & \sinh z = z + \frac{1}{3!}z^3 + \frac{1}{5!}z^5 + \dots \\
    & \tanh z = z - \frac{1}{3}z^3 + \frac{2}{15}z^5 + \dots \quad [|z| < \pi/2]
\end{align*}

\section{Integrals}
\begin{align*}
    & \int \frac{dx}{1 + x^2} = \tan^{-1} x \\
    & \int \frac{dx}{1 - x^2} = \tanh^{-1} x \\
    & \int \frac{dx}{\sqrt{1 - x^2}} = \sin^{-1} x \\
    & \int \frac{dx}{\sqrt{1 + x^2}} = \sinh^{-1} x \\
    & \int \tan x \, dx = - \ln \cos x \\
    & \int \tanh x \, dx = \ln \cosh x \\
    & \int \frac{dx}{x + x^2} = \ln \left( \frac{x}{1 + x} \right) \\
    & \int \frac{x \, dx}{1 + x^2} = \ln(1 + x^2) \\
    & \int \frac{dx}{\sqrt{x^2 - 1}} = \cosh^{-1} \, x \\
    & \int x \, dx \, \sqrt{1 + x^2} = \sqrt{1 + x^2} \\
    & \int \frac{dx}{x \sqrt{x^2 - 1}} = \cos^{-1}\left(\frac{1}{x}\right) \\
    & \int \frac{\sqrt{x} \, dx}{\sqrt{1 - x}} = \sin^{-1}(\sqrt{x}) - \sqrt{x(1 - x)} \\
    & \int \frac{dx}{(1 + x^2)^{3/2}} = \frac{x}{(1 + x^2)^{1/2}} \\
    & \int \ln(x) \, dx = x \ln(x) - x \\
\end{align*}






\end{document} % ----------Dokumentet er nu slut! ------------